\chapter{Parallelismo e Scalabilità}
Definiamo il parallelismo come l'\textbf{esecuzione simultanea} da parte di più unità di lavoro al fine di \textbf{ridurre il tempo di esecuzione o aumentare la dimensione del problema risolvibile.}
Facciamo una distinzione tra esecuzione in \textbf{contemporanea (concurrency)} ed esecuzione in \textbf{parallelo}:
\begin{itemize}
    \item \textbf{Concurrency}: La concurrency è la capacità di un sistema di gestire più compiti (task) nello stesso intervallo di tempo, ma \textbf{non necessariamente nello stesso istante.}
    \item \textbf{Parallelism}: Il parallelismo si verifica quando più compiti vengono eseguiti fisicamente nello stesso istante.
\end{itemize}
Adesso che abbiamo parlato dell'ottimizzazione, ci rendiamo facilmente conto di quanto il parallelismo sia ormai fondamentale nelle architetture odierne. Senza questo strumento sarebbe impossibile risolvere problemi complessi a cui lavoriamo oggi e, nonostante questo strumento, siamo ancora molto lontani dall'avere le prestazioni che vorremmo. 
\section{Modelli di parallelismo}
Parliamo dunque nel dettaglio di \textbf{modelli di parallelismo}. Questi definiscono:
\begin{itemize}
    \item Come il lavoro viene suddiviso;
    \item Come le unità di elaborazione interagiscono tra loro;
    \item Il modello di memoria da adottare;
    \item I costi di \textbf{comunicazione, sincronizzazione e gestione dei dati}.
\end{itemize}

Il modello ha, dunque, un'influenza riguardo al come dobbiamo progettare il nostro codice e di come questo possa scalare.
Possiamo dividere il parallelismo sulla base di \textbf{come si suddivide il lavoro}:
\begin{itemize}
    \item \textbf{Data Parallelism}: La stessa operazione è applicata a porzioni o \textbf{chunck} di dati diversi. Questo tipo di parallelismo è dominante in HPC scientifico.
    \item \textbf{Task Parallelism}: Compiti diversi vengono eseguiti in parallelo. Tipico in applicazioni eterogenee.
\end{itemize}
O sulla base di come è organizzata la memoria:
\begin{itemize}
    \item \textbf{Shared memory}: Secondo questa organizzazione abbiamo uno \textbf{spazio di indirizzamento unico}. Abbiamo quindi che i thread/processi guardano tutti la stessa porzione di memoria e, dunque, comunicano \textbf{implicitamente} proprio guardando la memoria. Questo modello di memoria è generalmente adottato sui singoli nodi poiché non è possibile collegare più di un certo numero di nodi alla stessa memoria per problemi di sovraffollamento e controlli di coerenza tra di essi.
    \item \textbf{Distribuited memory}: Ogni processo ha una memoria \textbf{privata} che, quindi, nel caso in cui avesse bisogno di comunicare con un altro processo, dovrà farlo mediante messaggi \textbf{espliciti}. Questo modello è generalmente usato nei cluster così ché ogni nodo del cluster abbia la sua memoria dato che questi sono indipendenti gli uni dagli altri, ciò è molto più scalabile.
    \item \textbf{Hibrid}: I sistemi HPC moderni sono spesso \textbf{ibridi} per cercare di prendere il meglio dall'uno e dall'altro. 
\end{itemize}
\subsection*{Modelli a Memoria Condivisa}
Come già detto, nel modello a memoria condivisa, tutte le unità di esecuzione condividono lo spazio degli indirizzi.
Le variabili sono dunque visibili da tutte le unità di esecuzione e la comunicazione avviene implicitamente tramite l'accesso a memoria condivisa.

Questo modello ha il vantaggio di avere una comunicazione intra-nodo estremamente rapida al costo, però, di dover \textbf{gestire i meccanismi di sincronizzazione} dovuti ai problemi di race condition ecc.
\subsection*{Modelli a Memoria Distribuita}
Nei modelli a memoria distribuita ogni processo ha la sua memoria privata e i dati che si scambiano i vari processi devono essere scambiati esplicitamente mediante messaggi o scambi di dati. Questo garantisce un \textbf{maggiore controllo sulla distribuzione dei dati} a patto di avere una maggiore complessità riguardo alla programmazione degli stessi. 

Questo modello è \textbf{altamente scalabile} e adatto a \textbf{sistemi multi-nodo}.
\section{Scalabilità}
La scalabilità è la capacità di un sistema parallelo di migliorare le prestazioni al crescere delle risorse utilizzate.

Distinguiamo tra \textbf{Strong Scaling} e \textbf{Weak Scaling}:
\begin{itemize}
    \item \textbf{Strong Scaling}: Questo scenario suppone di avere un problema di dimensione fissa e che, all'aumentare del numero di processi di lavoro, si misura una riduzione del tempo di risoluzione del problema.
    \item \textbf{Weak Scaling}: Qui, invece, abbiamo uno scenario con un problema che cresce all'aumentare dei processi. In questo caso, il nostro obiettivo sarà quello non di andare a diminuire il tempo di risoluzione del problema, bensì quello di farlo rimanere costante all'aumentare della dimensione del problema.
\end{itemize}

Sia $T(1)$ il tempo di risoluzione con un processo e $T(p)$ il tempo di risoluzione con $p$ processi, definiamo:
\begin{itemize}
    \item \textbf{Speedup}: Come il rapporto tra il tempo di esecuzione con un singolo processo e il tempo con $p$ processi, cioè:
    \[S(p) = \frac{T(1)}{T(p)}\]
    Questo va a misurare quanto il nostro problema migliora in termini di tempo all'aumentare dei processi che lo elaborano. Idealmente, vorremmo che $S(p) = p$.
    \item \textbf{Efficienza}: Questa è definita come il rapporto tra lo Speedup e p, cioè:
    \[E(p)=\frac{S(p)}{p}\]
    Questo va a misurare quanto il nostro problema scala bene all'aumentare dei processi che lo elaborano. Idealmente, vorremmo che $E(p)=1$.
\end{itemize}
In realtà, per quanto vorremmo che lo Speedup sia lineare, non è così poiché:
\begin{itemize}
    \item \textbf{Parte Seriale}: Una parte del programma potrebbe non essere parallelizzabile e, dunque, l'aumentare i processi non è una soluzione.
    \item \textbf{Comunicazione}: Inevitabilmente, a causa dell'uso di molti processi, questi hanno la necessità di comunicare e, ciò, introduce una latenza.
    \item \textbf{Sincronizzazione}: Poiché lavoriamo con sezioni critiche ecc. abbiamo la necessità di introdurre delle attese.
    \item \textbf{Overhead di gestione}: Dovuto allo scheduling o alla creazione dei thread/processi.
    \item \textbf{Load imbalance}: Cioè alcuni processi terminano prima di altri.
\end{itemize}
\subsection*{Legge di Amdahl}
Sia $\alpha$ la frazione parallelizzabile e $1-\alpha$ la parte seriale. Il tempo di esecuzione con $p$ processi è:
\[T(p)=T(1)* [(1-\alpha)+\frac{\alpha}{p}] \]
Lo Speedup è, invece:
\[S(p)=\frac{1}{(1-\alpha)+\frac{\alpha}{p}}\]
Il limite teorico per p che tende a infinito è:
\[\lim_{p\to \inf}S(p)= \frac{1}{1-\alpha}\]
Questo dipende dal fatto che, anche con infiniti processi, la parte seriale rimane un limite non superabile.
\subsection*{Legge di Gustafson}
La legge di Gustafson è applicabile nel caso di \textbf{weak scaling} in cui il problema cresce con il numero di processi ma la parte seriale rimane costante. Quella che, dunque, non sarà costante, sarà proprio la parte parallelizzabile che crescerà con p.
Lo Speedup diventerà:
\[S(p)=p-(1-\alpha)(p-1)\]
Se $\alpha$ fosse vicino a 1, cioè la parte parallelizzabile fosse quasi la totalità del programma, allora lo Speedup crescerebbe quasi linearmente. Per questa legge non abbiamo un limite come nel caso della legge di Amdahl.

Abbiamo un vantaggio nell'aggiungere processi quando:
\begin{itemize}
    \item La parte seriale è minima.
    \item Il problema è sufficientemente grande a tal punto che il calcolo domina la comunicazione.
    \item Il carico è ben bilanciato.
\end{itemize}